\documentclass[a4paper,11pt]{article}
\usepackage[utf8]{inputenc}
\usepackage[T1]{fontenc}
\usepackage[ngerman]{babel}
\usepackage{amsmath,amssymb}
\usepackage{xcolor}
\usepackage{colortbl}
\usepackage{array}
\usepackage{tikz}
\usepackage{enumitem}
\usepackage[margin=2cm]{geometry}
\definecolor{bgdark}{HTML}{1E1E1E}
\definecolor{fgtext}{HTML}{E6E6E6}
\definecolor{gridcol}{HTML}{4A4A4A}
\definecolor{accent}{HTML}{6CB6FF}
\definecolor{loescol}{HTML}{7FD18C}
\pagecolor{bgdark}
\color{fgtext}
\arrayrulecolor{fgtext}
\setlength{\parindent}{0pt}
\newcommand{\karo}[1]{\par\noindent\begin{tikzpicture}\draw[step=5mm,gridcol,very thin](0,0) grid (\linewidth,#1);\end{tikzpicture}\par\vspace{2mm}}
\newcommand{\aufgabe}[2]{\vspace{4mm}\par\noindent{\large\color{accent}\textbf{Aufgabe #1}}\hfill{\color{gridcol}\normalsize #2}\par\vspace{1mm}}
\newcommand{\titel}[1]{\begin{center}{\LARGE\color{accent}\textbf{Optimierungsverfahren}}\\[3pt]{\large #1}\end{center}\vspace{1mm}\noindent\textbf{Name:}\ \underline{\hspace{6cm}}\hfill\textbf{Datum:}\ \underline{\hspace{4cm}}\\[3pt]{\color{gridcol}\rule{\linewidth}{0.4pt}}\vspace{2mm}}

\begin{document}
\titel{Übungsaufgaben --- Blatt 03: Analytische \& restringierte Optimierung}

\aufgabe{1 --- Stationäre Punkte \& Hesse-Matrix}{mittel}
Gegeben $f(x_1,x_2)=x_1^3+x_2^3-3x_1-12x_2+7$.
\begin{enumerate}[label=(\alph*)]
  \item Bestimmen Sie alle stationären Punkte über $\nabla f=\vec 0$.
  \item Klassifizieren Sie jeden Punkt (Minimum / Maximum / Sattelpunkt) mit Hilfe der Hesse-Matrix.
\end{enumerate}
\karo{8cm}

\aufgabe{2 --- Definitheit}{leicht--mittel}
Bestimmen Sie für jede Matrix die Definitheit und den daraus folgenden Typ des stationären Punktes
(Min / Max / Sattel / unklar):
\[
H_1=\begin{pmatrix}4&1\\1&3\end{pmatrix},\quad
H_2=\begin{pmatrix}-2&1\\1&-3\end{pmatrix},\quad
H_3=\begin{pmatrix}1&2\\2&1\end{pmatrix},\quad
H_4=\begin{pmatrix}2&2\\2&2\end{pmatrix}.
\]
\karo{6cm}

\aufgabe{3 --- Funktion mit Mischterm}{mittel}
Gegeben $f(x_1,x_2)=x_1^2+2x_2^2+2x_1x_2-6x_1-8x_2$.
Bestimmen Sie den stationären Punkt und klassifizieren Sie ihn über die Hesse-Matrix.
Geben Sie auch den Funktionswert im Optimum an.
\karo{6.5cm}

\aufgabe{4 --- Gradientenverfahren}{mittel}
Gegeben $f(x_1,x_2)=x_1^2+2x_2^2$, Startpunkt $\vec x_0=(2,1)$, Schrittweite $\eta=0{,}1$.
Führen Sie \textbf{zwei} Schritte des Gradientenverfahrens $\vec x_{k+1}=\vec x_k-\eta\nabla f(\vec x_k)$ aus
und geben Sie jeweils den Funktionswert an.
\karo{6cm}

\aufgabe{5 --- Lagrange (Gleichungs-NB)}{mittel}
Lösen Sie mit der Lagrange-Methode:
\[
\min_{\vec x}\ f(\vec x)=x_1^2+x_2^2 \quad\text{u.d.N.}\quad h(\vec x):\ x_1+x_2-4=0.
\]
\karo{6cm}

\aufgabe{6 --- KKT aufstellen}{mittel}
Stellen Sie für das folgende Problem die \textbf{KKT-Bedingungen} auf (alle vier Bedingungen).
Das Problem muss \emph{nicht} gelöst werden.
\[
\min_{\vec x}\ f(\vec x)=x_1^2+x_2^2 \quad\text{u.d.N.}\quad g(\vec x):\ 1-x_1-x_2\le 0.
\]
\karo{6cm}

\aufgabe{7 --- KKT lösen}{mittel--schwer}
Lösen Sie das Problem aus Aufgabe 6 vollständig über die KKT-Bedingungen
(Fallunterscheidung Komplementärschlupf). Geben Sie $\vec x^\ast$, $b$ und $f(\vec x^\ast)$ an.
\karo{7cm}

\aufgabe{8 --- Quadratische Optimierung}{mittel}
Gegeben $f(x_1,x_2)=x_1^2+x_2^2+x_1x_2-3x_1$.
\begin{enumerate}[label=(\alph*)]
  \item Schreiben Sie $f$ in der Form $\tfrac12\,\vec x^\top Q\,\vec x+\vec c^\top\vec x$ (also $Q$ und $\vec c$ angeben).
  \item Ist $Q$ positiv definit (also $f$ konvex)?
  \item Bestimmen Sie das Minimum über $\nabla f=Q\vec x+\vec c=\vec 0$.
\end{enumerate}
\karo{6.5cm}

\aufgabe{9 --- Sequentielle quadratische Optimierung}{schwer}
Bei der SQO wird in jedem Schritt das Teilproblem
$\min_{\vec d}\ f(\vec x_i)+\nabla f(\vec x_i)^\top\vec d+\tfrac12\,\vec d^\top H\,\vec d$ gelöst.
Für $f$ aus Aufgabe 3 sei der Startpunkt $\vec x_0=(0,0)$.
\begin{enumerate}[label=(\alph*)]
  \item Geben Sie $\nabla f(\vec x_0)$ und $H$ an.
  \item Bestimmen Sie die Suchrichtung $\vec d$ aus $H\vec d=-\nabla f(\vec x_0)$ und den neuen Punkt $\vec x_1=\vec x_0+\vec d$.
  \item Was fällt auf?
\end{enumerate}
\karo{7cm}

\end{document}
